\section{Introduction}
\label{sec:introduction}

A precise understanding of parton
distributions~\cite{Forte:2010dt,Forte:2013wc,Ball:2015oha} (PDFs) has 
played a major role in the discovery of the Higgs boson and will be a 
key ingredient in searches for new physics at the LHC~\cite{Rojo:2015acz}. In recent years a new
generation of PDF sets~\cite{Ball:2014uwa,Dulat:2015mca,Harland-Lang:2014zoa,Alekhin:2017kpj,Abramowicz:2015mha,Jimenez-Delgado:2014twa,Accardi:2016qay} have 
been developed for use at the LHC Run II.
Some of these have been used in the construction of the PDF4LHC15 
combined sets,  
recommended for new physics searches and for the assessment of PDF uncertainties on
precision observables~\cite{Butterworth:2015oua}.  These PDF4LHC15 sets are
obtained by means of statistical combination of the three global
sets~\cite{Ball:2014uwa,Dulat:2015mca,Harland-Lang:2014zoa}: this is
justified by the improved level of agreement in the global
determinations, with differences between 
them largely consistent with statistical fluctuation. 

Despite these developments, there remains a need for improvements in the precision and
reliability of PDF determinations. Precision measurements at the LHC, 
such as in the search for new physics through Higgs
coupling measurements, will eventually require a systematic knowledge of PDFs at the percent 
level in order to fully exploit the LHC's potential. 
The NNPDF3.0 PDF set~\cite{Ball:2014uwa}, which is one of the sets entering the PDF4LHC15
combination, is unique in being a PDF set based on a methodology 
systematically validated by means of closure tests, which ensure the
statistical consistency of the procedure used to extract the PDFs from data.
The goal of this paper is to present NNPDF3.1, an update of the
NNPDF3.0 set, and  
a first step towards PDFs with percent-level uncertainties. Two directions of
progress are required in order to reach
this goal, the motivation for an update being accordingly twofold.

On the one hand, bringing the precision of PDFs down to the percent
level needs  a larger and more precise dataset, with correspondingly 
precise theoretical predictions. In the time since the release of NNPDF3.0, a significant 
number of new experimental measurements have become available.
From the Tevatron, we now have the final measurements of the 
$W$ boson asymmetries with the electron and muon final states 
based upon the complete
Run II dataset~\cite{Abazov:2013rja,D0:2014kma}. At the LHC, 
the ATLAS, CMS and LHCb experiments have released a wide variety
of measurements on inclusive jet production, gauge boson production and 
top production. Finally, the combined legacy measurements of DIS 
structure functions from HERA have also become available~\cite{Abramowicz:2015mha}.
In parallel with the experimental developments, an impressive number of new high-precision 
QCD calculations of hadron collider processes with direct sensitivity to PDFs have
recently been completed, enabling their use in the determination of PDFs 
at NNLO. These include differential distributions in top quark pair
production~\cite{Czakon:2015owf,Czakon:2016dgf}, the transverse 
momentum of the $Z$ and $W$ bosons~\cite{Boughezal:2015dva,Ridder:2015dxa},
and inclusive jet production~\cite{Currie:2013dwa,Currie:2016bfm}, for all
of which precision ATLAS and CMS datasets are available.

All of these new datasets and calculations have been incorporated into
NNPDF3.1. The inclusion of the new data presents new challenges. Given
the large datasets on which some of these measurements are based,
uncorrelated experimental uncertainties are often at the permille
level. Achieving a good fit then requires an unprecedented control of
both correlated systematics and of the numerical accuracy of
theoretical predictions. 

On the other hand, with uncertainties at the percent level, accuracy
issues related to theoretical uncertainties hitherto not included
in PDF determinations become relevant. Whereas the comprehensive inclusion of theoretical
uncertainties in PDF determination will require further study, we
have recently argued that a significant 
source of theoretical bias arises from the conventional assumption 
that charm is generated entirely perturbatively from gluons and light quarks. 
A methodology which allows for the inclusion
of a parametrized 
heavy quark PDFs within the FONLL matched general-mass variable
flavor number scheme has been developed~\cite{Ball:2015tna,Ball:2015dpa}, and 
implemented in an NNPDF PDF determination~\cite{Ball:2016neh}.
It was found that when the charm PDF is parametrized and determined
from the data  alongside the other PDFs, 
much of the uncertainty related to the value of the charm mass becomes part of
the standard PDF uncertainty, while any bias related to the assumption 
that the charm PDF is purely perturbative is eliminated~\cite{Ball:2016neh}.
In NNPDF3.1 charm is therefore parametrized as an independent PDF, 
in an equivalent manner to light quarks and the gluon. 
We will show that this leads to improvements in fit quality without an increase
in uncertainty, and that it stabilises the dependence of PDFs on the charm mass, 
all but removing it in the light quark PDFs. 

The NNPDF3.1 PDf sets are released at LO, NLO, and NNLO accuracy. For the first
time, all NLO and NNLO 
PDFs are delivered both as Hessian sets and as Monte Carlo
replicas, exploiting recent powerful methods for the construction of
optimal Hessian representations of 
PDFs~\cite{Carrazza:2015aoa}. Furthermore, and also 
for the first time, the default
PDF sets are provided as compressed Monte Carlo
sets~\cite{Carrazza:2015hva}. Therefore despite being presented
as sets of only 100 Monte Carlo replicas, they exhibit many of
the statistical properties of a much larger set, 
reducing observable computation time without loss of
information. A further improvement in computational efficiency can be
obtained by means of the {\tt SM-PDF} tool~\cite{Carrazza:2016htc}, which
allows for the selection of optimal subsets of Hessian eigenvectors
for the computation of uncertainties on specific processes or classes
of processes, and which is available as a web
interface~\cite{Carrazza:2016wte} now  also including the NNPDF3.1 sets.
A variety
of PDF sets based on subsets of data are also provided (as standard
100 replica Monte Carlo sets), which may be
useful for specific applications such as new physics searches, or measurements of
standard model parameters.

The outline of this paper is as follows.
%
First, in Sect.~\ref{sec:expdata} we discuss the
experimental aspects and the relevant theoretical
issues of the new datasets. 
We then turn in Sect.~\ref{sec:results} to a detailed
description of the baseline NNPDF3.1 PDF sets, with a specific discussion of 
the impact of methodological improvements, specifically the fact that
the charm PDF is now independently parametrized and determined like
all other PDFs. In
Sect.~\ref{sec:impactnewdata} we discuss the impact of the
new data by comparing PDF sets based upon various data subsets, and also
discuss PDF sets based on more conservative data subsets.
In Sect.~\ref{sec:pheno} we summarise the status of uncertainties on PDFs
and luminosities, and
specifically discuss the strange and charm content of the proton in
light of our results, and  present
first phenomenological studies at the LHC.
Finally, a summary of the PDFs being delivered in various formats  is
provided in 
Sect.~\ref{sec:conclusion}, together with links to repositories whence
more detailed sets  of plots may be downloaded.

%%%%%%%%%%%%%%%%%%%%%%%%%%%%%%%%%%%%%%%%%%%%%%%%%%%%%%%%%%%%%%%%%%%%%%%%%%%
%%%%%%%%%%%%%%%%%%%%%%%%%%%%%%%%%%%%%%%%%%%%%%%%%%%%%%%%%%%%%%%%%%%%%%%%%%%
